\subsection{P032 --- Graph neural network-based spatio-temporal indoor environment prediction and optimal control for central air-conditioning systems}
\label{paper:P032}
\textbf{Authors:} Jing Zhang, Fu Xiao, Ao Li, Tianyou Ma, Kan Xu, Hanbei Zhang, Rui Yan, Xing Fang, Yuanyang Li, Dan Wang\par
\textbf{Major Category:} AI and Machine Learning for Buildings\par
\textbf{Secondary Category:} Building Environment and IEQ, BIM and Semantic Information, BMS, BAS, BACS and Building Automation\par
\textbf{Keywords:} Air conditioning system, Graph neural network, Indoor environment, Machine learning, Optimal control, BIM, BAS, IoT\par
\paragraph{主な主張}
BAS, BIM, IoT由来のmulti-source dataをgraph structureとして統合し, GNN-RNNで室内温度分布を予測してonline optimal controlへ接続するspatio-temporal methodologyを実建物AHU-VAV systemで検証した. \cite{Zhang:2023GNNIndoorControl}
\paragraph{新規性}
単一計測値または平均値で室内環境を代表させる従来のcontrol formulationに対し, built environmentとHVAC systemのspatial relationshipをgraphとして明示し, distributed indoor environmentを予測と制御に直接利用した点に新規性がある.
\paragraph{位置付け}
BIM, BAS, IoTから構成したgraphをprediction modelの構造へ直接反映する例であり, semantic graphではないが, building informationをlearning architectureへ埋め込むspatio-temporal predictionの比較対象として位置付ける.
